% GENERATED FILE. Edit canonical lessons, then run npm run generate:books.
% canonical-source-hash: fnv1a64:ee13025302512bfc
% canonical-lessons: FR-C448-coiffeur, FR-C448-raser, FR-C448-opticien, FR-C448-horloger, FR-C448-bijoutier

\chapter{Coiffeur et opticien}
\label{ch:fr-coiffeur-et-opticien}
\hlchaptermodality{448}

\begin{chapteropening}
\textbf{By the end of this chapter:} \emph{I can use personal services: the hairdresser, shaving, the optician, the watchmaker and the jeweller.}

Complete the last lesson of chapter 448: I can use personal services: the hairdresser, shaving, the optician, the watchmaker and the jeweller.
\end{chapteropening}

\section[le coiffeur]{le coiffeur — the hairdresser, the barber}

\label{lesson:FR-C448-coiffeur}



\begin{quote}
Before the new one: say the French for to dye, then the French for to alter.
\end{quote}

\subsection*{You'll want to know: le coiffeur}
\textbf{le coiffeur} — \textquotedblleft{}the hairdresser, the barber\textquotedblright{}.

\begin{quote}
Je vais chez le coiffeur une fois par mois.
\end{quote}

\begin{quote}
\emph{I go to the hairdresser's once a month.}
\end{quote}

Feminine \emph{la coiffeuse}. \emph{Aller chez le coiffeur}.

\begin{cousinweb}
\emph{coiffeur} is built on \emph{coiffer} 'to do someone's hair' (EtymDB), from \emph{coiffe} 'cap', which is usually traced to Late Latin \emph{cofia}, of Germanic origin.
\end{cousinweb}

\begin{grammarlens}[title={What you've built}]
One more word for telling, explaining and describing at B1.
\end{grammarlens}

\subsection*{Your turn}
\begin{itemize}
\raggedright
  \item \emph{Say it:} \emph{le coiffeur}
  \item \emph{Say it:} \emph{le coiffeur}, once more
  \item \emph{Say it:} a sentence of your own with \emph{le coiffeur}
  \item \emph{Recall:} say the French for to dye, then the French for to alter, then say \emph{le coiffeur} again
\end{itemize}

\begin{tcolorbox}[breakable,colback=teal!4,colframe=teal!35!black,title={Before you move on}]
What does le coiffeur mean? (\textquotedblleft{}The hairdresser, the barber\textquotedblright{}.) Say it once more.
\end{tcolorbox}
\section[raser]{raser — to shave (often se raser); to raze}

\label{lesson:FR-C448-raser}



\begin{quote}
Before the new one: say the French for to alter, then the French for the hairdresser.
\end{quote}

\subsection*{You'll want to know: raser}
\textbf{raser} — \textquotedblleft{}to shave (often se raser); to raze\textquotedblright{}.

\begin{quote}
Le coiffeur lui a rasé la barbe.
\end{quote}

\begin{quote}
\emph{The barber shaved off his beard.}
\end{quote}

\emph{Se raser} — to shave; informal \emph{ça me rase} — that bores me.

\begin{cousinweb}
\emph{raser} is inherited from Vulgar Latin \emph{rāsāre}, from Latin \emph{rāsus} 'scraped', the past participle of \emph{rādere} 'to scrape' (EtymDB).
\end{cousinweb}

\begin{grammarlens}[title={What you've built}]
One more word for telling, explaining and describing at B1.
\end{grammarlens}

\subsection*{Your turn}
\begin{itemize}
\raggedright
  \item \emph{Say it:} \emph{raser}
  \item \emph{Say it:} \emph{raser}, once more
  \item \emph{Say it:} a sentence of your own with \emph{raser}
  \item \emph{Recall:} say the French for to alter, then the French for the hairdresser, then say \emph{raser} again
\end{itemize}

\begin{tcolorbox}[breakable,colback=teal!4,colframe=teal!35!black,title={Before you move on}]
What does raser mean? (\textquotedblleft{}To shave (often se raser); to raze\textquotedblright{}.) Say it once more.
\end{tcolorbox}
\section[l'opticien]{l'opticien — the optician}

\label{lesson:FR-C448-opticien}



\begin{quote}
Before the new one: say the French for the hairdresser, then the French for to shave.
\end{quote}

\subsection*{You'll want to know: l'opticien}
\textbf{l'opticien} — \textquotedblleft{}the optician\textquotedblright{}.

\begin{quote}
L'opticien a réparé mes lunettes tout de suite.
\end{quote}

\begin{quote}
\emph{The optician fixed my glasses straight away.}
\end{quote}

A masculine noun; feminine \emph{l'opticienne}.

\begin{cousinweb}
\emph{opticien} is built on \emph{optique} 'optics', a learned word from Latin \emph{opticus}, from Greek \emph{optikós} 'of seeing' (EtymDB). EtymDB lists \emph{opticien} itself without a source.
\end{cousinweb}

\begin{grammarlens}[title={What you've built}]
One more word for telling, explaining and describing at B1.
\end{grammarlens}

\subsection*{Your turn}
\begin{itemize}
\raggedright
  \item \emph{Say it:} \emph{l'opticien}
  \item \emph{Say it:} \emph{l'opticien}, once more
  \item \emph{Say it:} a sentence of your own with \emph{l'opticien}
  \item \emph{Recall:} say the French for the hairdresser, then the French for to shave, then say \emph{l'opticien} again
\end{itemize}

\begin{tcolorbox}[breakable,colback=teal!4,colframe=teal!35!black,title={Before you move on}]
What does l'opticien mean? (\textquotedblleft{}The optician\textquotedblright{}.) Say it once more.
\end{tcolorbox}
\section[l'horloger]{l'horloger — the watchmaker, the clockmaker}

\label{lesson:FR-C448-horloger}



\begin{quote}
Before the new one: say the French for to shave, then the French for the optician.
\end{quote}

\subsection*{You'll want to know: l'horloger}
\textbf{l'horloger} — \textquotedblleft{}the watchmaker, the clockmaker\textquotedblright{}.

\begin{quote}
L'horloger a changé la pile de ma montre.
\end{quote}

\begin{quote}
\emph{The watchmaker changed the battery in my watch.}
\end{quote}

A masculine noun. \emph{L'horlogerie} — the watchmaker's shop.

\begin{cousinweb}
\emph{horloger} is built on \emph{horloge} 'clock', which comes from Latin \emph{hōrologium}, from Greek \emph{hōrológion} 'thing that tells the hour' (Lewis \& Short). EtymDB does not list it.
\end{cousinweb}

\begin{grammarlens}[title={What you've built}]
One more word for telling, explaining and describing at B1.
\end{grammarlens}

\subsection*{Your turn}
\begin{itemize}
\raggedright
  \item \emph{Say it:} \emph{l'horloger}
  \item \emph{Say it:} \emph{l'horloger}, once more
  \item \emph{Say it:} a sentence of your own with \emph{l'horloger}
  \item \emph{Recall:} say the French for to shave, then the French for the optician, then say \emph{l'horloger} again
\end{itemize}

\begin{tcolorbox}[breakable,colback=teal!4,colframe=teal!35!black,title={Before you move on}]
What does l'horloger mean? (\textquotedblleft{}The watchmaker, the clockmaker\textquotedblright{}.) Say it once more.
\end{tcolorbox}
\section[le bijoutier]{le bijoutier — the jeweller}

\label{lesson:FR-C448-bijoutier}



\begin{quote}
Before the new one: say the French for the optician, then the French for the watchmaker.
\end{quote}

\subsection*{You'll want to know: le bijoutier}
\textbf{le bijoutier} — \textquotedblleft{}the jeweller\textquotedblright{}.

\begin{quote}
Il a acheté une bague chez le bijoutier.
\end{quote}

\begin{quote}
\emph{He bought a ring at the jeweller's.}
\end{quote}

\emph{La bijouterie} — the jeweller's shop.

\begin{cousinweb}
\emph{bijoutier} is built on \emph{bijou} 'jewel' (EtymDB lists it), which was borrowed from Breton \emph{bizou} 'ring', from \emph{biz} 'finger'.
\end{cousinweb}

\begin{grammarlens}[title={What you've built}]
One more word for telling, explaining and describing at B1.
\end{grammarlens}

\subsection*{Your turn}
\begin{itemize}
\raggedright
  \item \emph{Say it:} \emph{le bijoutier}
  \item \emph{Say it:} \emph{le bijoutier}, once more
  \item \emph{Say it:} a sentence of your own with \emph{le bijoutier}
  \item \emph{Recall:} say the French for the optician, then the French for the watchmaker, then say \emph{le bijoutier} again
\end{itemize}

\begin{tcolorbox}[breakable,colback=teal!4,colframe=teal!35!black,title={Before you move on}]
What does le bijoutier mean? (\textquotedblleft{}The jeweller\textquotedblright{}.) Say it once more.
\end{tcolorbox}
